\section{Classi, Ereditarietà e Generatori}

\section{Il Diabolico Dualismo: Proto vs Prototype}
Una delle fonti di confusione maggiori in JavaScript è la distinzione tra `__proto__` e `prototype`.

\begin{definition}{Distinzione Fondamentale}
    \begin{itemize}
        \item \textbf{\texttt{\_\_proto\_\_}}: È una proprietà presente in \textbf{TUTTI} gli oggetti (funzioni incluse). Rappresenta il legame alla catena dei prototipi (da chi eredito?).
        \item \textbf{\texttt{prototype}}: È una proprietà presente \textbf{SOLO} nelle funzioni. È l'oggetto che verrà assegnato come `__proto__` alle istanze create invocando quella funzione con `new`.
    \end{itemize}
\end{definition}

\begin{concept}{Relazione di Identità}
Data una funzione costruttore `F` e un oggetto `o` creato con `var o = new F()`, vale sempre:
$$ \texttt{F.prototype} \equiv \texttt{o.\_\_proto\_\_} $$
Nota bene: `F.__proto__` esiste, ma è il prototipo della funzione stessa (ovvero `Function.prototype`), non c'entra con gli oggetti creati da F.
\end{concept}

\section{Classi ed Espressioni Classe}
Le classi (introdotte in ES6) sono "zucchero sintattico" sopra il modello a prototipi esistente.

\subsection{Traduzione Dietro le Quinte}
Cosa succede realmente quando dichiariamo una classe?

\begin{jsbox}[title={Classe ES6}]
class Persona {
    constructor(n, e) {
        this.nome = n;
        this.eta = e;
    }
    compleanno() { this.eta++; }
}
\end{jsbox}

\begin{jsbox}[title={Equivalente ES5 (Pseudocodice)}]
function Persona(nome, eta) {
    this.nome = nome;
    this.eta = eta;
    // this.__proto__ = Persona.prototype (implicito nel new)
}
Persona.prototype = {
    constructor: Persona,
    compleanno: function() { this.eta++; }
};
\end{jsbox}

\subsection{Espressioni Classe}
Come le funzioni, le classi possono essere anonime e assegnate a variabili:
\begin{jsbox}
var Persona = class {
    constructor(n, e) { ... }
};
\end{jsbox}

\section{Ereditarietà Avanzata}

\subsection{Meccanismo di Implementazione}
L'ereditarietà `class Studente extends Persona` non fa altro che concatenare i prototipi:
\begin{enumerate}
    \item \textbf{Istanze:} `Studente.prototype.__proto__ === Persona.prototype`. Questo permette alle istanze di Studente di accedere ai metodi di Persona.
    \item \textbf{Statici:} `Studente.__proto__ === Persona`. Questo permette alla classe Studente di ereditare i metodi statici di Persona.
\end{enumerate}

\subsection{Overriding e `super`}
Una sottoclasse può ridefinire (override) i metodi.
\begin{itemize}
    \item \textbf{super.metodo()}: Invoca il metodo dal prototipo del genitore.
    \item \textbf{super() nel costruttore}: Obbligatorio prima di usare `this`.
\end{itemize}

\begin{algo}{Esecuzione di new Sottoclasse()}
Quando si esegue `new Studente(...)`:
\begin{enumerate}
    \item Viene creato un oggetto vuoto.
    \item La catena dei costruttori viene risalita fino alla radice (`Persona`).
    \item Viene eseguito il costruttore di `Persona` (inizializza proprietà base).
    \item Viene eseguito il costruttore di `Studente` (aggiunge proprietà specifiche).
    \item L'oggetto completo viene restituito.
\end{enumerate}
\end{algo}

\section{Membri di Classe (ES2022+)}

\subsection{Campi Pubblici, Privati e Statici}
JavaScript moderno permette definizioni dichiarative.

\begin{jsbox}[title={Esempio Completo Liceale}]
class Liceale extends Studente {
    // Campo Statico (della classe)
    static quanti = 0;

    // Campo Privato (visibile solo dentro la classe)
    #iscritto = false;

    // Campo Pubblico (inizializzato su ogni istanza)
    materie = ["Mate", "Italiano"];

    constructor(nome) {
        super(nome);
        Liceale.quanti++;
    }

    unipi() { return this.#iscritto; } // OK
    
    // Accesso esterno a #iscritto -> Errore SyntaxError
}
\end{jsbox}

\section{Getter e Setter (Accessors)}
Permettono di intercettare la lettura e scrittura di proprietà.

\begin{jsbox}
let o = {
    get x() { return 1; },      // Lettura o.x restituisce sempre 1
    set x(v) { console.log(v); } // Scrittura o.x = 5 logga 5
};
\end{jsbox}

\begin{algo}{Strategia di Ricerca Proprietà}
Quando si accede a `obj.key`:
\begin{enumerate}
    \item Se la chiave esiste nell'oggetto:
    \begin{itemize}
        \item Se è un \textbf{getter}, invoca la funzione e restituisce il risultato.
        \item Se è una proprietà base, restituisce il valore.
    \end{itemize}
    \item Se non esiste, cerca nel prototipo (ricorsivamente).
    \item Se non trovata: \texttt{undefined}.
\end{enumerate}
Similmente per la scrittura (con i \textbf{setter}). I setter vengono ereditati!
\end{algo}

\section{Generatori}
I generatori sono funzioni speciali che possono sospendere l'esecuzione e riprenderla, mantenendo lo stato locale.

\subsection{Sintassi e Funzionamento}
\begin{itemize}
    \item Dichiarazione: \texttt{function* nome()} o \texttt{*metodo()}.
    \item \textbf{\texttt{yield expr}}: Restituisce \texttt{expr} al chiamante e "congela" la funzione.
    \item \textbf{\texttt{next(val)}}: Riprende l'esecuzione. Se si passa \texttt{val}, questo diventa il risultato dell'espressione \texttt{yield} che era in attesa.
\end{itemize}

\begin{jsbox}[title={Esempio Range}]
function* range(a, b) {
    let i = a;
    while(i < b) yield i++;
}
// Uso
for (var n of range(4, 8)) { console.log(n); } // 4, 5, 6, 7
\end{jsbox}

\subsection{Metodi Avanzati dei Generatori}
Oltre a `next()`, l'oggetto generatore possiede:
\begin{itemize}
    \item \textbf{\texttt{return(v)}}: Termina forzatamente il generatore (come se ci fosse un `return v` nel codice).
    \item \textbf{\texttt{throw(e)}}: Lancia un'eccezione all'interno del generatore (nel punto dello `yield`).
\end{itemize}

\subsection{Generatori Ricorsivi (`yield*`)}
Per delegare la generazione a un altro iterabile, si usa `yield*`.
\begin{jsbox}
function* merge(g1, g2) {
    yield* g1; // Emette tutti i valori di g1
    yield* g2; // Poi tutti i valori di g2
}
\end{jsbox}

\subsection{Esempio Complesso: La Poesia}
Un generatore che estrae parole da un testo.
\begin{jsbox}
class Poesia {
    #testo;
    constructor(t) { this.#testo = t; }

    *parole() {
        var i = 0;
        while (true) {
            let f = this.#testo.indexOf(" ", i); // Cerca spazio
            if (f >= 0) {
                yield this.#testo.slice(i, f); // Restituisce parola
                i = f + 1; // Avanza indice
            } else {
                yield this.#testo.slice(i); // Ultima parola
                break; // Fine
            }
        }
    }
}
var p = new Poesia("La vispa Teresa...");
console.log([...p.parole()]); // Array di parole
\end{jsbox}

\section{Esercizi: Code (Queue)}
Implementazione completa richiesta dal laboratorio.

\subsection{Classe Base: Queue}
\begin{jsbox}
class Queue {
    #array = []; // Array privato per lo storage

    put(x) {
        // Aggiunge in fondo
        this.#array[this.#array.length] = x;
    }

    get() {
        // Rimuove e restituisce dalla testa
        let x = this.#array[0];
        this.#array = this.#array.slice(1);
        return x;
    }

    size() { return this.#array.length; }
    
    empty() { return this.#array.length === 0; }
    
    toString() {
        return `Queue with ${this.size()} elements`;
    }
}
\end{jsbox}

\subsection{Sottoclasse: LimitedQueue}
Eredita da Queue e aggiunge un controllo sulla capacità massima.

\begin{jsbox}
class LimitedQueue extends Queue {
    #max_size; // Capacità massima

    constructor(max_size) {
        super(); // Inizializza this tramite Queue
        this.#max_size = max_size;
    }

    // Metodo aggiuntivo
    full() {
        return this.size() == this.#max_size;
    }

    // Override di put
    put(x) {
        if (this.full()) {
            console.log('Queue is full!!');
        } else {
            super.put(x); // Delega alla superclasse
        }
    }

    // Override di toString con super
    toString() {
        return super.toString() + ` of ${this.#max_size}`;
    }
}
\end{jsbox}
